For the validation of  properties checked in previous chapters
and for demonstrating the feasibility of 
the counterexample-guided control approach, we complement the formal analysis 
with an experimental simulation platform. The platform is realized in 
the Gazebo environment, which provides a physics-based simulation of 
the modular robot and its interaction with voxel-like blocks. The robot 
is controlled through ROS2 nodes, enabling modular and distributed control 
strategies. By comparing the outcomes of formal verification with simulation 
results, we aim to assess both the soundness and the practical utility of 
model-checking-based methods for modular self-replicating robots.

\section{Validation of the counterexample based robot control}

To validate the counterexample based approach, introduced in 
chapter \ref{ch:Counterexample} we will integrate the counterexample generation 
pipeline into ROS2-Gazebo simulation environment.
We choose robot blocks configuration described in Chapter \ref{ch:robotSimulation}, 
robot initial position and target point with coordinates $[50,50,10]$. 
These initial parameters reflected in the lines 1-3 of the
 Listing \ref{lst:validCounterInit} lines.

\begin{lstlisting}[language=Python, caption={Function goto\_XYZ() of RobotController class}, 
    label={lst:validCounterInit}]
angle_vars = ["block1.yaw", "block2.pitch", "block3.pitch", "block4.pitch", "block5.yaw"]
sequence = ["base", "yaw", "pitch", "pitch", "pitch", "yaw"]
lengths = [10, 20, 30, 30, 10, 20]
robot = Robot(sequence, lengths)
generator = RobotNusmvGenerator(robot, output_smv_path="./smv/")
generator.set_target_point(50, 50, 10)
generator.generate_from_template("project_robot")
result = Counterexample.run_nusmv_file("./smv/robot.smv")
counterexample = Counterexample(result, len(sequence))
angles = counterexample.get_angles(angle_vars)
counterexample.plot_trace(-1, show_plot=True)
\end{lstlisting}
Trace derived from the counterexample consist out of 9 steps and 
shown in Figure \ref{fig:validCounterexamplePlot}
\begin{figure}[h!]
    \centering
    \includegraphics[height=10cm]{validCounterexamplePlot.png} 
    \caption{Mean time of model checking with number of blocks in robot}
    \label{fig:validCounterexamplePlot}
\end{figure}

Resulting joint angles derived from the trace presented as a list steps,
where each step is a list of control angles for each joint:
\[ [0,\ \ \ \ 0,\ \ \ \ 0,\ \ \ \ 0,\ \ \ \ 0] \]
\[ [0, -10, -10, -10, -10] \]
\[ [0, -20, -20, -20, -20] \]
\[ [0, -30, -30, -30, -30] \]
\[ [0, -30, -40, -40, -40] \]
\[ [0, -30, -50, -50, -50] \]
\[ [0, -30, -60, -50, -60] \]
\[ [0, -30, -70, -50, -60] \]
\[ [0, -30, -80, -50, -60] \]

We will publish these angles as a sequence of commands for robot joints 
using \texttt{rotate\_joints()} function in RobotController class, described in
Chapter \ref{ch:robotSimulation}.
Before simulation we have to adjust position and rotation of the gazebo coordinate system 
and coordinate system recieved from counterexample.
The simulation result can be seen on the Fugure \ref{fig:validCounterSimulation}

\begin{figure}[h!]
    \centering
    % \includegraphics[height=4cm]{counterSimulation/0.png} 
    \includegraphics[height=4.5cm]{counterSimulation/1.png} 
    \includegraphics[height=4.5cm]{counterSimulation/2.png} 
    \includegraphics[height=4.5cm]{counterSimulation/3.png} 
    \includegraphics[height=4.5cm]{counterSimulation/4.png} 
    \includegraphics[height=4.5cm]{counterSimulation/5.png} 
    \includegraphics[height=4.5cm]{counterSimulation/6.png} 
    \includegraphics[height=4.5cm]{counterSimulation/7.png} 
    \includegraphics[height=4.5cm]{counterSimulation/8.png} 
    \includegraphics[height=4.5cm]{counterSimulation/9.png} 
    \caption{Simulation steps to achive point [5,5,1] from counterexample}
    \label{fig:validCounterSimulation}
\end{figure}

For considered target point the distance from robot end to target is around 0.16.

Total discrepancy between the target point and the center of the 
end effector consists of the angle discretization error in NuSMV and 
the geometric offset of the blocks in the simulation model. 
Maximal discrepancy during simulation of 10 different reachable points 
was 27\% of length unit in the simulation (size of the voxel).

\section{Validation of the robot configuration checking}

\section{Complexity evaluation}
\textcolor{red}{TODO: Add comparison of a time to preform an assembly action}\\
\textcolor{red}{TODO: Add graph of a model checking time for different lengths}\\

We can assess the complexity of the configuration checking with sequential 
model checking, using a target point with coordinates: x=5, y=5, z=10. 
The changing of mean time required for checking relative to a number of blocks in 
robot shown in the Figure \ref{fig:ConfigCheckingPointTimeSize}

\begin{figure}[h!]
    \centering
    \includegraphics[height=10cm]{Time_size_plot.png} 
    \caption{Mean time of model checking with number of blocks in robot}
    \label{fig:ConfigCheckingPointTimeSize}
\end{figure}
